%% AION-1 talk — beamer/moloch
%% Build: latexmk -lualatex presentation.tex   (or use the top-level Makefile)

\documentclass[10pt]{beamer}

\usetheme{moloch}

% ── fonts (Fira Sans via fontspec; requires LuaLaTeX or XeLaTeX) ──────────────
\usepackage{fontspec}
\setsansfont{Fira Sans}[
  BoldFont = {Fira Sans Bold},
  ItalicFont = {Fira Sans Italic},
]
\setmonofont{Fira Mono}

% ── packages ──────────────────────────────────────────────────────────────────
\usepackage{graphicx}
\usepackage{booktabs}
\usepackage{amsmath,amssymb}
\usepackage{xcolor}
\usepackage{tikz}

% ── bibliography ──────────────────────────────────────────────────────────────
\usepackage[backend=biber, style=authoryear, maxnames=2]{biblatex}
\addbibresource{bibliography/presentation-bibliography.bib}
% To also draw on the AION-1 paper bibliography (available after make fetch):
% \addbibresource{../resources/aion-1-paper/mmoma.bib}

% ── graphicspath ──────────────────────────────────────────────────────────────
\graphicspath{%
  {../resources/papers/aion-1-paper/source/figures/}%
  {../resources/papers/melchior2023spender/source/figures/}%
  {../resources/papers/ginolfi2024moons/source/figures/}%
  {../resources/papers/duev2019realbogus/source/}%
  {../resources/}%
  {../resources/images/}%
  {../resources/papers/}%
}

% ── moloch customisation ──────────────────────────────────────────────────────
\molochset{progressbar=frametitle}
\setbeamertemplate{frame numbering}[fraction]

% ── metadata ──────────────────────────────────────────────────────────────────
\title{AION-1: Omnimodal Foundation Model for Astronomical Sciences}
\author{Tom Hehir}
\institute{Institute of Astronomy, University of Cambridge}
\date{2 March 2026}

% ── convenience macros ────────────────────────────────────────────────────────
\newcommand{\aion}{AION-1}

% ══════════════════════════════════════════════════════════════════════════════
\begin{document}

\maketitle

% ── §2 Abstract ───────────────────────────────────────────────────────────────
\begin{frame}{Abstract}
  \centering
  \includegraphics[height=0.82\textheight, trim={0 207pt 0 0}, clip]{titlepage.pdf}
  \par\vspace{0.5em}
  \href{https://arxiv.org/abs/2510.17960}{%
    \includegraphics[height=1.6em]{badge-arxiv.pdf}%
  }%
  \hspace{0.5em}%
  \href{https://neurips.cc/virtual/2025/loc/san-diego/poster/119776}{%
    \includegraphics[height=1.6em]{badge-neurips.pdf}%
  }
\end{frame}

% ── §3 Team ───────────────────────────────────────────────────────────────────
\begin{frame}{The Team}
  \centering
  \includegraphics[width=\textwidth]{team-collage.png}
  \par\vspace{0.5em}
  \href{https://polymathic-ai.org}{%
    \includegraphics[height=1.2em]{polymathic-logo.pdf}%
  }
\end{frame}

% ── §3 Status Quo ─────────────────────────────────────────────────────────────
\begin{frame}{Status Quo: Task- and Data-Specific Models}
  \begin{columns}[T]

    \begin{column}{0.32\textwidth}
      \centering
      \includegraphics[height=0.42\textheight]{sketch.pdf}
      \par\smallskip
      {\small\textbf{Physical inductive bias}}\\[0.3em]
      {\small
        Architecture encodes redshift and instrument response explicitly
        \parencite{melchior2023spender};
        or via the training objective
        \parencite{korber2023pinion}.
      }
    \end{column}

    \begin{column}{0.32\textwidth}
      \centering
      \includegraphics[height=0.42\textheight]{nn-sketch.png}
      \par\smallskip
      {\small\textbf{Data-specific regression}}\\[0.3em]
      {\small
        Multi-task net with probabilistic outputs for joint redshift and
        galaxy property estimation on simulated MOONS spectra
        \parencite{ginolfi2024moons}.
        Bayesian SED emulation: \parencite{alsing2020speculator}.
      }
    \end{column}

    \begin{column}{0.32\textwidth}
      \centering
      \includegraphics[height=0.42\textheight]{fig-braai3.png}
      \par\smallskip
      {\small\textbf{Survey-specific classification}}\\[0.3em]
      {\small
        VGG6 on ZTF science/reference/difference triplets for real-bogus
        discrimination \parencite{duev2019realbogus};
        CNNs for strong lens finding \parencite{metcalf2019lensfinding}.
      }
    \end{column}

  \end{columns}
\end{frame}

\end{document}
